% Owner: Ky. Write LAST, after Section 3 numbers are final.
\begin{abstract}
\noindent
\owner{Ky (write last)}
% [Ky]: draft (about 230 words) following the thesis proposed below; revise once the title is decided.
Clustering methods differ in what they take a cluster to be, and real data break those ideas in
different ways. We compare K-means, K-means++, DBSCAN and HDBSCAN on eight datasets, each chosen to test
one data property: curved shapes, long belts of uneven density, density contrast within a city, rare
classes, high dimensions, classes of unequal spread, and data with no gaps between groups. We ask where
each method fails, whether the scores available without labels reveal those failures, and whether known
remedies repair them. Each method failed where the data broke its idea of a cluster. K-means cut the
curved classes of make\_moons (ARI 0.48 against 0.98 for DBSCAN); one DBSCAN radius clustered 90\% of
earthquakes in the dense south-west Pacific but 20\% in the Americas; and in 7 to 50 dimensions the
density methods fell to ARI 0.03 to 0.05. The silhouette did not reveal these failures: on five of six
labelled datasets the method with the top silhouette was not the one that matched the labels best, and
on Ecoli it ranked the four methods almost in reverse (Spearman $\rho=-0.80$). Rerunning each result on
fake data caught methods that found nothing real, but not ones that found the wrong structure. Surface
distance, HDBSCAN and standardisation each fixed a specific failure at a cost; none produced clusters in
data with no gaps. The choice of method should follow from the structure of the data, and every score
needs a check against fake data.


\decide{Thesis, proposed 7 Oct (Wei Yew), to match the research question in Section~\ref{sec:intro}:
each method fails where the data break its idea of a cluster (RQ1); the silhouette does not reveal
these failures, a fake-data check reveals only some, and the rest need checks aimed at a known problem
(RQ2); known remedies fix some failures at a cost, and none fixes data with no gaps (RQ3).}
\end{abstract}
